\section{AI Native 架构: 七层 KOC 决策智能 OS}

\subsection{为什么不是流水线, 而是 OS}

市面上的 KOC 内容工具普遍是\textbf{ Agent 流水线} (\textit{输入选题 → 写文案 → 配图 → 发布}), 这种结构有三个根本缺陷:

\begin{itemize}
    \item \textbf{决策不可审}: KOC 问 \textit{「为什么是这个话题」} 时, 流水线只能给最终答案, 看不到否决了哪些备选, 用了什么证据。
    \item \textbf{风格不持久}: 每次 prompt 重新喂一遍人设, 风格在跨会话间会漂移, 一段时间后 KOC 自己都觉得 \textit{「不像我了」}。
    \item \textbf{能力不组合}: 流水线的能力是焊死的, 想要 \textit{「先扫描全网信号, 再用我的人设过滤, 然后做 7 天战役规划」}这种组合, 必须重写整条流水线。
\end{itemize}

\begin{insight}
Ripple 2.0 的核心设计:把 KOC 的内容决策建模成一个 \textbf{操作系统 (OS)}, 而不是单一应用。
OS 提供原语 (内存管理 / 调度 / 文件系统 / 进程隔离), 应用通过组合原语完成任务。
Ripple 2.0 提供的原语是 \textbf{40+ Tool + 5 层记忆 + Persona Vector + Replay Graph}, 用户的一句话被规划器拆解为 DAG, 然后调度原语去解决。
\end{insight}

\subsection{七层架构总览}

\begin{figure}[ht]
\centering
\begin{tikzpicture}[
    layer/.style={rectangle, rounded corners=4pt, draw=ripple-primary, fill=ripple-light,
                  minimum width=14cm, minimum height=1.0cm, align=center, font=\small},
    arrow/.style={->, >=Stealth, thick, draw=ripple-accent}
]
    \node[layer] (l7) at (0, 0) {\textbf{L7 Distribution} \quad Web SSE \textbar{} MCP Skill \textbar{} Tauri 桌面 \textbar{} CLI};
    \node[layer] (l6) at (0, -1.2) {\textbf{L6 Persistence} \quad SQLite (relational) \textbar{} Vector Store \textbar{} Markdown \textbar{} Event Log};
    \node[layer] (l5) at (0, -2.4) {\textbf{L5 Decision (原创)} \quad SimPredictor \textbar{} Risk-Reward \textbar{} TrendChain \textbar{} Cohort \textbar{} Translator \textbar{} Campaign};
    \node[layer] (l4) at (0, -3.6) {\textbf{L4 Reasoning} \quad Planner \textbar{} Critic \textbar{} Forum Debate \textbar{} Verifier \textbar{} Reflector \textbar{} Replay Recorder};
    \node[layer] (l3) at (0, -4.8) {\textbf{L3 Action} \quad Tool Registry (40+) \textbar{} Subagent Dispatcher \textbar{} Permission Sandbox};
    \node[layer] (l2) at (0, -6.0) {\textbf{L2 Cognition} \quad LLM Router \textbar{} Persona Vector \textbar{} Skill Library \textbar{} Knowledge Base};
    \node[layer] (l1) at (0, -7.2) {\textbf{L1 Sensing} \quad 7 数据源接入 \textbar{} 变点检测 (CUSUM/MAD) \textbar{} 事件知识图谱 \textbar{} Embedding};

    \draw[arrow] (l7) -- (l4);
    \draw[arrow] (l4) -- (l5);
    \draw[arrow] (l4) -- (l3);
    \draw[arrow] (l3) -- (l1);
    \draw[arrow] (l4) -- (l2);
    \draw[arrow] (l2) -- (l6);
    \draw[arrow] (l5) -- (l6);
    \draw[arrow] (l1) -- (l6);
\end{tikzpicture}
\caption{Ripple 2.0 七层架构总览}
\end{figure}

\subsection{L1 Sensing 感知层}

\textbf{职责}: 把外部世界 (社交平台 / 预测市场 / 资本流) 的原始信号变成结构化数据。

\textbf{关键能力}:
\begin{itemize}
    \item \textbf{7 数据源并行扫描}: Polymarket / Manifold / HackerNews 直连; 微博 / 抖音 / 百度 / B 站经第三方聚合 (\texttt{xxapi.cn})
    \item \textbf{变点检测}: CUSUM 累积和检测 + MAD-zscore 鲁棒统计 (适应噪声)
    \item \textbf{事件知识图谱}: 实体-话题-平台-时间四元组
    \item \textbf{Embedding}: 字符 trigram hash (256 维) + 可替换为 sentence-transformers
\end{itemize}

\textbf{输出对象}: \texttt{HotItem (source, title, rank, raw\_value, normalized, url, metadata)} 与 \texttt{Citation (url, title, source\_type, retrieved\_at, snippet, confidence)}。

\subsection{L2 Cognition 认知层}

\textbf{职责}: AI 思考所需的工作内存与长期知识。

\textbf{核心模块}:
\begin{itemize}
    \item \textbf{LLM Router}: 多 Provider 抽象 (默认 MiMo, Fallback MiniMax → DeepSeek → Hunyuan), 真流式 SSE, 调用观测
    \item \textbf{Persona Vector}: 256 维 + 10 维可解释指标 (详见 §\ref{sec:cap-persona})
    \item \textbf{Skill Library}: Markdown 加载 + 自进化 (仿 Claude Code)
    \item \textbf{Knowledge Base}: 向量检索 + 关键词过滤
\end{itemize}

\subsection{L3 Action 行动层}

\textbf{职责}: 把能力做成可组合的工具, 由上层 Reasoning 调度。

\textbf{工具注册表 (Tool Registry)}: 一个工具的定义包含

\begin{lstlisting}[language=Python, caption={Tool 注册示例}]
@register_tool(
    name="oracle.scan",
    description="并行扫描 7 数据源, 返回 HotItem 矩阵",
    permission_level=PermissionLevel.NETWORK,
    input_model=OracleScanInput,
    output_model=OracleScanOutput,
    cost_estimate_ms=15000,
    tags=["sensing", "oracle"],
)
async def oracle_scan(input, ctx, bus): ...
\end{lstlisting}

每次工具调用都被 \textbf{ReplayRecorder} 自动记录: 输入哈希 + 输出摘要 + 耗时 + Merkle 链, 支持任意节点重放。

\textbf{权限分级}: \texttt{READ < NETWORK < GENERATE < WRITE < DESTRUCTIVE}, 高权限工具调用需用户明示授权 (仿 Claude Code 权限分层)。

\subsection{L4 Reasoning 推理层}

\textbf{职责}: 决定 \textit{要做什么}, 是 OS 的 \textbf{kernel scheduler}。

\textbf{TAOR 认知循环}:

\begin{figure}[ht]
\centering
\begin{tikzpicture}[
    node distance=0.6cm,
    box/.style={rectangle, rounded corners, draw=ripple-primary, fill=ripple-light,
                minimum width=2.0cm, minimum height=0.8cm, font=\small\bfseries},
    arrow/.style={->, >=Stealth, thick}
]
    \node[box] (sense) {Perceive};
    \node[box, right=of sense] (think) {Think};
    \node[box, right=of think] (plan) {Plan};
    \node[box, right=of plan] (act) {Act};
    \node[box, right=of act] (obs) {Observe};
    \node[box, right=of obs] (refl) {Reflect};
    \node[box, right=of refl] (upd) {Update};

    \draw[arrow] (sense) -- (think);
    \draw[arrow] (think) -- (plan);
    \draw[arrow] (plan) -- (act);
    \draw[arrow] (act) -- (obs);
    \draw[arrow] (obs) -- node[above, font=\tiny]{未达成} (think.south east);
    \draw[arrow] (obs) -- (refl);
    \draw[arrow] (refl) -- (upd);
    \draw[arrow, dashed] (upd) to[bend right=18] node[above, font=\tiny]{更新 Skill/Persona/Memory} (think);
\end{tikzpicture}
\caption{TAOR + Reflection 认知循环}
\end{figure}

\textbf{核心模块}:
\begin{itemize}
    \item \textbf{Adaptive Planner}: 用 LLM 把用户的一句话拆为 \texttt{TaskDAG}, 选择需要哪些 Skill 和 Tool
    \item \textbf{Critic}: 对 Planner 的决策提反对意见, 生成 counter-arguments
    \item \textbf{Forum Debate}: 多人格 (KOC / 法务 / 营销) 辩论, 共识收敛
    \item \textbf{Verifier}: 多源交叉验证, 计算一致性 score
    \item \textbf{Reflector}: 任务结束后反思, 更新 Skill / Persona
    \item \textbf{Replay Recorder}: 记录每一步 ReplayNode + Merkle 摘要
\end{itemize}

\subsection{L5 Decision 决策层 (Ripple 原创)}

这是 Ripple 区别于通用 Agent 框架的关键, 提供 6 大决策能力 (详细见 §\ref{sec:capabilities-10}):

\begin{tabularx}{\textwidth}{>{\bfseries\color{ripple-primary}}l X}
    \toprule
    SimPredictor & 虚拟受众团 (5 cohort) 对多版本内容评分, 自动选最优 \\
    \midrule
    Risk-Reward Scorer & 选题双头打分 (爆款潜力 vs 翻车风险) + 监管敏感度 + 推荐动作 \\
    \midrule
    Trend Chain Analyzer & 资本→国际→中文→大众 四阶段因果链推理 \\
    \midrule
    Cohort Affinity & 内容对人群的吸引力预测 + 破圈建议 \\
    \midrule
    Cross-Platform Translator & 一条源内容 → 多平台 (小红书/视频号/B站/抖音/微博/公众号) 改写 \\
    \midrule
    Campaign Planner & 7 天战役图 (流量结构约束 + MIP/贪心调度) \\
    \bottomrule
\end{tabularx}

\subsection{L6 Persistence 持久化层}

\textbf{多存储统一抽象}:

\begin{itemize}
    \item \textbf{SQLite} (\texttt{users / sessions / persona\_vectors / memory\_entries / replay\_nodes / audit\_log / campaigns / topics / event\_log})
    \item \textbf{Vector Store}: 用于语义检索 (chromadb / lance, 当前用 Hash trigram 占位)
    \item \textbf{Markdown 文件}: \texttt{\textasciitilde/.ripple/memory/<user>/} 用户级记忆 (仿 Claude Code MEMORY.md)
    \item \textbf{Event Log}: append-only 事件溯源, 是 Replay 的底层
\end{itemize}

\textbf{多层记忆系统}:

\begin{figure}[ht]
\centering
\begin{tabularx}{\textwidth}{>{\bfseries\color{ripple-primary}}l X X}
    \toprule
    层 & 内容 & 生命周期 \\
    \midrule
    User & 用户偏好 / 边界 / 订阅级别 & 永久 \\
    Persona & 人设向量 + 风格规则 + 禁忌 & 长期 (有 EMA 增量更新) \\
    Project & 当前营销项目 / 系列内容 / 历史决策 & 项目期 \\
    Session & 当前对话上下文 & 会话内 \\
    Long-term & 跨项目知识 / 历史 Replay 摘要 & 长期 (后台 Dream 巩固) \\
    \bottomrule
\end{tabularx}
\end{figure}

\textbf{后台 Dream 巩固} (仿 Claude Code 后台压缩): 定时任务巡检日志, 把活跃 Session 压缩为 long-term 摘要, 合并矛盾, 重写指针索引。

\subsection{L7 Distribution 分发层}

\textbf{三端一致}:

\begin{itemize}
    \item \textbf{Web App} (Next.js 风格静态 + SSE): \url{http://120.55.247.6/chat}
    \item \textbf{MCP Skill} (\texttt{npx @ripple/mcp}): 在 Cursor / Claude Code 里调用 Ripple 工具集
    \item \textbf{Tauri 桌面端} (Mac/Win): 系统托盘 + 全局快捷键 + 离线 SQLite
    \item \textbf{CLI / TUI}: \texttt{ripple ask "<query>"} 命令行直接调用
\end{itemize}

\subsection{核心数据流}

一次完整 Aha Demo 的数据流:

\begin{figure}[ht]
\centering
\begin{tikzpicture}[
    node distance=0.4cm and 0.5cm,
    box/.style={rectangle, rounded corners, draw=ripple-primary, fill=ripple-light,
                minimum width=2.6cm, minimum height=0.7cm, font=\footnotesize, align=center},
    arrow/.style={->, >=Stealth, thick, draw=ripple-accent}
]
    \node[box] (q) {用户一句话};
    \node[box, right=of q] (p) {Planner};
    \node[box, right=of p] (mem) {Memory Recall};
    \node[box, below=of mem] (oracle) {Oracle Scan};
    \node[box, left=of oracle] (rr) {Risk-Reward};
    \node[box, left=of rr] (sim) {SimPredictor};
    \node[box, below=of sim] (camp) {Campaign Planner};
    \node[box, right=of camp] (crit) {Self-Critique};
    \node[box, right=of crit] (cite) {Citation Enforce};
    \node[box, below=of cite] (final) {最终结论 + Replay};

    \draw[arrow] (q) -- (p);
    \draw[arrow] (p) -- (mem);
    \draw[arrow] (mem) -- (oracle);
    \draw[arrow] (oracle) -- (rr);
    \draw[arrow] (rr) -- (sim);
    \draw[arrow] (sim) -- (camp);
    \draw[arrow] (camp) -- (crit);
    \draw[arrow] (crit) -- (cite);
    \draw[arrow] (cite) -- (final);
\end{tikzpicture}
\caption{一次完整 Aha Demo 的数据流}
\end{figure}

每个箭头都是流式事件 (\texttt{SSE}), 前端实时展示思考过程, 这就是评委看到的 \textit{「Gemini Deep Research 风格」} 体验。
